% !TEX root = ./main.tex

%%%%%%%%%%%%%%%%%%%%%%%%%%%%%%%%%%%%%%%%%%%%%%%%%%%%%%
\section{改行}
%%%%%%%%%%%%%%%%%%%%%%%%%%%%%%%%%%%%%%%%%%%%%%%%%%%%%%
\TeX での改行を正しく理解しましょう．
普通に改行しても生成される文章は改行されませんが，
あああああああああああああああああああああああああああああああああああ

\TeX ソース内で空行を空けて改行すると，段落が変わったと認識され自動的に字下げされます．
いいいいいいいいいいいいいいいいいいいいいいいいいいいいいいいいいいい
いいいいいいいいいいいいいいいいいいいいいいいいいいいいいいいいいいい．

生成される文章で空行を開けたい場合は

\vspace*{1zh}
\noindent このようにします．
{\verb+\noindent+}を指示すると，字下げしなくなりますが，通常は指示せず文章を構成しましょう．

%%%%%%%%%%%%%%%%%%%%%%%%%%%%%%%%%%%%%%%%%%%%%%%%%%%%%%
\section{数式}\label{sec:math}
数式は文章の一部とされるので，以下に示すように本文中に差し込んでも良い．
%%%%%%%%%%%%%%%%%%%%%%%%%%%%%%%%%%%%%%%%%%%%%%%%%%%%%%
\subsection{数式の使い方}
%%%%%%%%%%%%%%%%%%%%%%%%%%%%%%%%%%%%%%%%%%%%%%%%%%%%%%
\subsubsection{数式の使い方その 1}
%%%%%%%%%%%%%%%%%%%%%%%%%%%%%%%%%%%%%%%%%%%%%%%%%%%%%%
数式を文章中に入れるには，「\verb+\+( \textbf{式} \verb+\)+」または「\verb+$+ \textbf{式} \verb+$+」で囲んで
\(
	\dot{x} = {A}{x} + {B}{u}
\) または $\dot{x} = Ax + Bu$ のようにします．また，「\verb+\[+ \textbf{式} \verb+\]+」で囲んで
\[
	\dot{x} = Ax + Bu
\]
とすれば，改行後に数式が挿入されます．式番号は「\verb+\begin{equation}+ \textbf{式} \verb+\end{equation}+」で囲んで
\begin{equation}
	\dot{x} = Ax + Bu
\label{eq:state}
\end{equation}
とすれば自動的につきますし，\eqref{eq:state}式 のように，式番号を参照することもできます．
\verb+\ref+と\verb+\eqref+の違いも知っておきましょう．

%%%%%%%%%%%%%%%%%%%%%%%%%%%%%%%%%%%%%%%%%%%%%%%%%%%%%%
\subsubsection{数式の使い方その 2}
%%%%%%%%%%%%%%%%%%%%%%%%%%%%%%%%%%%%%%%%%%%%%%%%%%%%%%
「{\verb+\begin{align}+ \textbf{式} \verb+\end{align}+}」で囲んで
\begin{align}
	E\dot{x} &= Ax + Bu \\
	y &= Cx + Du
\end{align}
とすれば，「\verb+&+ \textbf{そろえたい部分}」で (``\,=\,'') をそろえることができます．\verb+\notag+で
\begin{align}
	E\dot{x} &= Ax + Bu \notag \\
	y &= Cx + Du
\end{align}
のように式番号をはずすこともできます．

%%%%%%%%%%%%%%%%%%%%%%%%%%%%%%%%%%%%%%%%%%%%%%%%%%%%%%
\subsection{数式中の書体}
%%%%%%%%%%%%%%%%%%%%%%%%%%%%%%%%%%%%%%%%%%%%%%%%%%%%%%
\subsubsection{斜体の解除}
数式中の変数は自動的に斜体になりますが，定数や添え字など正体で表現したい場合は
\[
	{v}(t) = {k}_\mathrm{P}e(t) 
			+ {k}_\mathrm{I}\int_{0}^{t}{e(\tau)d\tau}
	\ \ \Longleftrightarrow\ \ 
	{v}(s) = \left({k}_\mathrm{P} + \frac{{k}_\mathrm{I}}{s}\right){e}(s)
\]
のように\verb+\mathrm+を使います．

三角関数や対数などの関数は
\[
	\tan\theta = \frac{\sin\theta}{\cos\theta}
\]
のように\verb+\sin+などの命令を使いましょう．これらを斜体にしてはいけません．

%%%%%%%%%%%%%%%%%%%%%%%%%%%%%%%%%%%%%%%%%%%%%%%%%%%%%%
\subsubsection{ボールドイタリック}
「\verb+\bm{ボールドイタリックにしたい部分}+」により数式中の文字をボールドイタリックにすることができます．
\begin{equation}
	\bm{\dot{x}} = \bm{Ax} + \bm{Bu}
\end{equation}

%%%%%%%%%%%%%%%%%%%%%%%%%%%%%%%%%%%%%%%%%%%%%%%%%%%%%%
\subsection{ディスプレイ形式の分数}
%%%%%%%%%%%%%%%%%%%%%%%%%%%%%%%%%%%%%%%%%%%%%%%%%%%%%%
「\verb+\frac{分子}{分母}+」で表現された分数を行列中などに用いると，
\begin{equation}
	\left[ \begin{array}{cc}
		a & \frac{c}{b} \\
		\frac{e}{b + c} & f
	\end{array} \right]
\end{equation}
のように分数が小さくなってしまいます．通常の大きさにしたい場合には，「\verb+\frac{分子}{分母}+」の代わりに「\verb+\dfrac{分子}{分母}+」を使用して下さい．「\verb+\dfrac{分子}{分母}+」を用いると，次式のようになります．
\begin{equation}
	\left[ \begin{array}{cc}
		a & \dfrac{c}{b} \\
		\dfrac{e}{b + c} & f
	\end{array} \right]
\end{equation}
